\documentclass[aspectratio=169]{beamer}
\usetheme{metropolis}
\usepackage{appendixnumberbeamer}
\usepackage{booktabs, hyperref}

\input{mgmt675-style}

\usepackage{tikz}
\usetikzlibrary{shapes.geometric, arrows.meta, positioning, calc}

% Title info
\subtitle{MGMT 675: Generative AI for Finance}
\title{Module 5: Automating Financial Workflows}
\author{Kerry Back}
\date{}

\begin{document}

\maketitle

% ========================================
% SECTION 1: THE DASHBOARD TRAP
% ========================================
\section{The Dashboard Trap}

\begin{frame}{Dashboards Answer Yesterday's Questions}
Organizations spend millions building dashboards. When a new question arises, the cycle restarts: requirements $\rightarrow$ design $\rightarrow$ build $\rightarrow$ deploy.
\vspace{0.3cm}
\begin{columns}[T]
\begin{column}{0.5\textwidth}
\begin{shadedbox}[title=\textbf{The Typical Dashboard Lifecycle}]
\begin{enumerate}\small
  \item Business user requests a report
  \item Team builds queries, charts, and deploys
  \item User asks a follow-up --- \alert{back to step 1}
\end{enumerate}
\end{shadedbox}
\end{column}
\begin{column}{0.5\textwidth}
\begin{shadedbox}[title=\textbf{The Cost}]
\begin{itemize}\small
  \item \textbf{Time}: Weeks to months per dashboard
  \item \textbf{Money}: BI licenses, engineering hours, maintenance
  \item \textbf{Rigidity}: Fixed views of fixed data
\end{itemize}
\end{shadedbox}
\vspace{0.2cm}
Gartner: only \textbf{20\%} of analytic insights deliver business outcomes.
\end{column}
\end{columns}
\end{frame}

\begin{frame}{The Fundamental Problem}
Dashboards answer \textbf{pre-defined questions}. But the most valuable analysis comes from \textbf{ad-hoc questions} that arise in the moment.
\vspace{0.3cm}
\begin{itemize}
  \item ``What happened to margins in the Southeast last quarter?''
  \item ``Show me our top 10 customers by growth rate, excluding one-time orders''
  \item ``Compare Q3 headcount vs.\ budget by department, and flag anyone over 110\%''
\end{itemize}
\vspace{0.3cm}
These are simple questions. Getting answers shouldn't require a development cycle.
\end{frame}

% ========================================
% SECTION 2: NATURAL LANGUAGE AS INTERFACE
% ========================================
\section{Natural Language as the Query Interface}

\begin{frame}{The Shift: From Dashboards to Conversations}
\begin{center}
\begin{tabular}{@{} l l l @{}}
\toprule
& \textbf{Traditional Dashboard} & \textbf{Natural Language AI} \\
\midrule
Query method & Click filters, select dates & Ask in plain English \\
Time to answer & Minutes to weeks & Seconds \\
Follow-up questions & New dashboard request & Next sentence \\
Visualization & Pre-built charts & Generated on demand \\
Who can use it & Trained users & Anyone \\
Cost per question & High (amortized) & Near zero (marginal) \\
\bottomrule
\end{tabular}
\end{center}
\vspace{0.3cm}
\centering
The dashboard was a \textit{workaround} for the fact that databases don't speak English. Now they do.
\end{frame}

\begin{frame}{What This Looks Like in Practice}
\begin{columns}[T]
\begin{column}{0.5\textwidth}
\begin{shadedbox}[title=\textbf{The Conversation}]
\begin{itemize}\small
  \item \textit{``Show me monthly revenue by product line for 2025''}
  \item AI: writes SQL, produces grouped bar chart
  \item \textit{``Break out Enterprise by region''}
  \item AI: refines query, updates chart
\end{itemize}
\end{shadedbox}
\end{column}
\begin{column}{0.5\textwidth}
\textbf{What the User Needed to Know}
\begin{itemize}\small
  \item What questions to ask
  \item Whether the answers make sense
  \item \alert{Nothing else}
\end{itemize}
\vspace{0.1cm}
\textbf{Behind the Scenes}
\begin{itemize}\small
  \item 4 different SQL queries written
  \item 3 visualizations produced
  \item Derived metrics calculated
\end{itemize}
\end{column}
\end{columns}
\end{frame}

\begin{frame}{The Database Agent Pattern}
The most powerful dashboard replacement: an AI agent connected to your database.
\vspace{0.3cm}
\begin{columns}[T]
\begin{column}{0.5\textwidth}
\begin{shadedbox}[title=\textbf{How to Build It}]
\begin{enumerate}\small
  \item Connect database via MCP or file upload
  \item Give AI the schema: table names, columns, relationships
  \item Describe the business context
  \item Start asking questions
\end{enumerate}
\end{shadedbox}
\end{column}
\begin{column}{0.5\textwidth}
\begin{shadedbox}[title=\textbf{What the Agent Can Do}]
\begin{itemize}\small
  \item Write and execute SQL queries
  \item Compute derived metrics (growth rates, ratios)
  \item Generate charts and export to Excel or PowerPoint
\end{itemize}
\end{shadedbox}
\end{column}
\end{columns}
\end{frame}

\begin{frame}{Replacing Dashboards: FP\&A and Treasury}
\begin{columns}[T]
\begin{column}{0.5\textwidth}
\begin{shadedbox}[title=\textbf{FP\&A}]
\begin{itemize}\small
  \item \textit{``Budget vs.\ actual for Q3, decompose the variance into volume, price, and cost drivers''}
  \item \textit{``Rolling 12-month revenue forecast with confidence bands''}
\end{itemize}
\end{shadedbox}
\end{column}
\begin{column}{0.5\textwidth}
\begin{shadedbox}[title=\textbf{Treasury \& Risk}]
\begin{itemize}\small
  \item \textit{``Cash position over the next 90 days using AR/AP forecasts''}
  \item \textit{``VaR by desk for the last 30 days, flag breaches''}
\end{itemize}
\end{shadedbox}
\end{column}
\end{columns}
\vspace{0.3cm}
Each of these would take a BI team \textbf{days to build}. With a database agent, they take \alert{seconds}.
\end{frame}

\begin{frame}{Replacing Dashboards: Portfolio and Executive}
\begin{columns}[T]
\begin{column}{0.5\textwidth}
\begin{shadedbox}[title=\textbf{Portfolio Management}]
\begin{itemize}\small
  \item \textit{``Sector allocation vs.\ benchmark with active weights''}
  \item \textit{``Performance attribution --- allocation vs.\ selection''}
\end{itemize}
\end{shadedbox}
\end{column}
\begin{column}{0.5\textwidth}
\begin{shadedbox}[title=\textbf{Executive Reporting}]
\begin{itemize}\small
  \item \textit{``One-page executive summary with KPIs and trends''}
  \item \textit{``Board deck from this quarter's financials --- 5 slides max''}
\end{itemize}
\end{shadedbox}
\end{column}
\end{columns}
\vspace{0.3cm}
The AI doesn't replace the analyst's judgment --- it replaces the \alert{mechanical work} of pulling data and building charts.
\end{frame}

\begin{frame}{The Weekly Ops Review --- Before and After}
\begin{columns}[T]
\begin{column}{0.5\textwidth}
\begin{shadedbox}[title=\textbf{Before: The Dashboard Era}]
\begin{enumerate}\small
  \item Data team pulls exports and builds slides (6+ hrs)
  \item Manager reviews and revises (3 hrs)
  \item VP asks a question --- ``We'll get back to you next week''
\end{enumerate}
\vspace{0.2cm}
\centering\small\textbf{Total: 9+ hours per week}
\end{shadedbox}
\end{column}
\begin{column}{0.5\textwidth}
\begin{shadedbox}[title=\textbf{After: Natural Language AI}]
\begin{enumerate}\small
  \item AI agent generates ops review (3 min)
  \item Manager reviews and iterates in chat
  \item VP asks a question --- \alert{AI answers in 10 seconds}
\end{enumerate}
\vspace{0.2cm}
\centering\small\textbf{Total: 15 minutes + live Q\&A}
\end{shadedbox}
\end{column}
\end{columns}
\end{frame}

% ========================================
% SECTION 3: SKILLS
% ========================================
\section{Skills: Encoding Expertise in Reusable Instructions}

\begin{frame}{It's All Just Text}
Every AI customization --- skills, slash commands, custom instructions --- works the same way: \alert{additional text is added to the prompt} that the model sees.
\vspace{0.3cm}
\begin{itemize}
  \item \textbf{Skills} = instructions for \textit{how} (loaded when relevant)
  \item \textbf{Slash commands} = triggers that say \textit{do} (typing \texttt{/name} loads that skill)
  \item Same pattern everywhere: Custom GPTs, Copilot instructions, \texttt{.cursorrules}, Gems
\end{itemize}
\end{frame}

\begin{frame}{What is a Skill?}
A \textbf{skill} is a set of instructions that specializes a general-purpose AI for a specific domain or task. In Claude, it's a folder with a markdown file.
\vspace{0.3cm}
\begin{columns}[T]
\begin{column}{0.5\textwidth}
\begin{shadedbox}[title=\textbf{A Skill Provides}]
\begin{itemize}\small
  \item System prompt (domain knowledge)
  \item Workflow instructions
  \item Best practices and constraints
\end{itemize}
\end{shadedbox}
\end{column}
\begin{column}{0.5\textwidth}
\begin{shadedbox}[title=\textbf{The AI Platform Provides}]
\begin{itemize}\small
  \item The LLM (Opus or Sonnet)
  \item Agent control loop and tools
  \item Code execution sandbox
\end{itemize}
\end{shadedbox}
\end{column}
\end{columns}
\vspace{0.3cm}
Skills let you create \alert{specialized agents without writing agent logic}.
\end{frame}

\begin{frame}[fragile]{Anatomy of a Skill}
\begin{columns}[T]
\begin{column}{0.45\textwidth}
\textbf{Skill Folder Structure}\small
\begin{verbatim}
skills/
  xlsx/
    SKILL.md     <- Main file
    scripts/
      recalc.py
    references/
      schema.md
\end{verbatim}
\end{column}
\begin{column}{0.55\textwidth}
\textbf{SKILL.md Structure}\small
\begin{verbatim}
---
name: xlsx
description: "Excel file..."
---

# Requirements for Outputs
- Zero formula errors
- Use formulas, not hardcodes

# Workflows
1. Choose pandas or openpyxl
2. Create/modify file
3. Recalculate formulas
\end{verbatim}
\end{column}
\end{columns}
\end{frame}

\begin{frame}{Where Skills Live}
\begin{itemize}
  \item \textbf{Project skills}: \texttt{.claude/skills/} in your project folder
  \item \textbf{User skills}: \texttt{\textasciitilde/.claude/skills/} (shared across projects)
  \item Skills also work in \textbf{Claude.ai} (upload ZIP), \textbf{Cowork} (install as plugin), and \textbf{VS Code}
\end{itemize}
\vspace{0.3cm}
\begin{center}
\small
\begin{tabular}{@{} l l l @{}}
\toprule
\textbf{Format} & \textbf{How to Install} & \textbf{How to Invoke} \\
\midrule
Claude.ai (web) & Upload ZIP in Settings & Automatic or \texttt{/name} \\
Code tab / CLI & Place in \texttt{.claude/skills/} & Automatic or \texttt{/name} \\
VS Code extension & Place in \texttt{.claude/skills/} & Automatic or \texttt{/name} \\
\bottomrule
\end{tabular}
\end{center}
\end{frame}

% ========================================
% SECTION 4: THE CONSISTENCY PROBLEM
% ========================================
\section{The Consistency Problem}

\begin{frame}{Why Skills Matter for Finance}
In Module 3, you built DCF models and portfolio optimizations in ad-hoc conversations. What happens when you start a new conversation?
\vspace{0.3cm}
\begin{columns}[T]
\begin{column}{0.5\textwidth}
\begin{shadedbox}[title=\textbf{The Problem}]
\begin{itemize}\small
  \item One conversation uses 10\% WACC, another uses 8\%
  \item One assumes 3\% terminal growth, another 2.5\%
  \item Portfolio optimization: 5-year history vs.\ 3-year
\end{itemize}
\end{shadedbox}
\end{column}
\begin{column}{0.5\textwidth}
\begin{shadedbox}[title=\textbf{The Fix: Skills}]
\begin{itemize}\small
  \item A DCF skill specifies: which assumptions, how to estimate them, what output format
  \item A portfolio skill specifies: data source, estimation window, constraints
  \item \alert{Every conversation follows the same playbook}
\end{itemize}
\end{shadedbox}
\end{column}
\end{columns}
\vspace{0.3cm}
\alert{No audit trail, no reproducibility --- unless you encode the methodology in a skill.}
\vspace{0.2cm}
\begin{center}
This is how AI moves from \textit{personal productivity tool} to \textit{institutional capability}.
\end{center}
\end{frame}

\begin{frame}{Finance Skill Examples: Analysis}
\begin{columns}[T]
\begin{column}{0.5\textwidth}
\begin{shadedbox}[title=\textbf{DCF Skill}]
\begin{itemize}\small
  \item Sales growth from trailing 5-year CAGR
  \item Sector median EV/EBITDA for exit multiple
  \item Two-way sensitivity table (WACC vs.\ growth)
\end{itemize}
\end{shadedbox}
\end{column}
\begin{column}{0.5\textwidth}
\begin{shadedbox}[title=\textbf{Investment Memo Skill}]
\begin{itemize}\small
  \item Standard structure: exec summary, financials, risks, recommendation
  \item Required data points and formatting
  \item Every memo follows the same template
\end{itemize}
\end{shadedbox}
\end{column}
\end{columns}
\vspace{0.3cm}
\alert{Every analyst produces comparable, auditable outputs.}
\end{frame}

\begin{frame}{Finance Skill Examples: Data}
\begin{columns}[T]
\begin{column}{0.5\textwidth}
\begin{shadedbox}[title=\textbf{Portfolio Optimization Skill}]
\begin{itemize}\small
  \item Data: Rice Data Portal, trailing 60 months
  \item Constraints: no short sales, max 25\%
  \item Output: weights, frontier plot, statistics
\end{itemize}
\end{shadedbox}
\end{column}
\begin{column}{0.5\textwidth}
\begin{shadedbox}[title=\textbf{Rice Data Portal Skill}]
\begin{itemize}\small
  \item Table schemas (SEP, DAILY, SF1, TICKERS)
  \item Query patterns and API access via \texttt{.env}
  \item Claude Code fetches data directly
\end{itemize}
\end{shadedbox}
\end{column}
\end{columns}
\vspace{0.3cm}
\alert{Consistent methodology across quarterly rebalancing --- no manual intervention.}
\end{frame}

\begin{frame}[fragile]{Installing the Rice Data Portal Skill}
In Module 1 you queried the Rice Data Portal through its web interface.  Now let's give Claude Code direct access.

\vspace{0.3cm}
\begin{columns}[T]
\begin{column}{0.5\textwidth}
\begin{shadedbox}[title=\textbf{Quick Install}]
\begin{itemize}\small
  \item Download \href{https://mgmt675.kerryback.com/files/rice-data-query.zip}{rice-data-query.zip} from the course site
  \item Tell Claude Code: ``Install this skill''
  \item Claude will unzip it to \texttt{.claude/skills/} and add a summary to your \texttt{CLAUDE.md}
\end{itemize}
\end{shadedbox}
\end{column}
\begin{column}{0.5\textwidth}
\begin{shadedbox}[title=\textbf{What It Provides}]
\begin{itemize}\small
  \item Connection setup and \texttt{.env} token handling
  \item Table schemas (SEP, DAILY, SF1, TICKERS)
  \item SQL query patterns and precomputed variable checker
\end{itemize}
\end{shadedbox}
\end{column}
\end{columns}

\vspace{0.3cm}
\textbf{Try it:} ``Get Apple's monthly returns for 2020--2024 and plot cumulative performance.''
\end{frame}

% ========================================
% SECTION 5: VARIANCE ANALYSIS
% ========================================
\section{Variance Analysis with the Finance Plugin}

\begin{frame}{What is Variance Analysis?}
\textbf{Variance analysis} compares \alert{budgeted} figures to \alert{actual} results and decomposes the differences into actionable drivers. It is the core analytical task in FP\&A.
\vspace{0.3cm}
\begin{columns}[T]
\begin{column}{0.5\textwidth}
\textbf{Why It Matters}
\begin{itemize}\small
  \item Every public company does it quarterly
  \item Boards and investors ask ``why did we miss?''
  \item Drives re-forecasting and capital allocation
\end{itemize}
\end{column}
\begin{column}{0.5\textwidth}
\textbf{Key Decompositions}
\begin{itemize}\small
  \item \textbf{Revenue}: volume vs.\ price vs.\ mix
  \item \textbf{COGS}: volume vs.\ unit cost
  \item \textbf{SG\&A}: headcount vs.\ rate vs.\ discretionary
\end{itemize}
\end{column}
\end{columns}
\vspace{0.3cm}
Revenue variance = \alert{Volume effect} + \alert{Price effect} + \alert{Mix effect}
\end{frame}

\begin{frame}[shrink=5]{Variance Analysis: Example Data}
\textbf{Scenario:} You are an FP\&A analyst. Q1 actuals just closed. The CEO wants to know why operating income missed budget by \$300K.
\vspace{0.2cm}
\begin{center}
\scriptsize
\begin{tabular}{@{} l r r r @{}}
\toprule
& \textbf{Budget} & \textbf{Actual} & \textbf{\$ Variance} \\
\midrule
\textbf{Revenue} & & & \\
\quad Units sold & 100,000 & 95,000 & \\
\quad Avg price & \$50.00 & \$51.00 & \\
\quad \textit{Total Revenue} & \textit{\$5,000,000} & \textit{\$4,845,000} & \textit{(\$155,000)} \\
\midrule
\textbf{COGS} & & & \\
\quad Unit cost & \$30.00 & \$32.00 & \\
\quad \textit{Total COGS} & \textit{\$3,000,000} & \textit{\$3,040,000} & \textit{(\$40,000)} \\
\midrule
\textit{Gross Profit} & \textit{\$2,000,000} & \textit{\$1,805,000} & \textit{(\$195,000)} \\
\midrule
\textit{Total SG\&A} & \textit{\$1,450,000} & \textit{\$1,555,000} & \textit{(\$105,000)} \\
\midrule
\textbf{Operating Income} & \textbf{\$550,000} & \textbf{\$250,000} & \textbf{(\$300,000)} \\
\bottomrule
\end{tabular}
\end{center}
\end{frame}

\begin{frame}[shrink=5]{Using the Finance Plugin}
With the finance plugin installed, one command replaces 2--4 hours of FP\&A work:
\vspace{0.2cm}

\textbf{Prompt}\\
\texttt{/finance:variance-analysis}\\[0.3em]
\small ``Here is our Q1 budget vs.\ actuals spreadsheet. Decompose the \$300K operating income miss into volume, price, cost, and discretionary spending drivers.  Produce a waterfall chart and a summary memo for the CFO.''
\normalsize
\vspace{0.2cm}

\textbf{What the Plugin Does}
\begin{enumerate}\small
  \item \textbf{Skill activates}: domain knowledge about variance decomposition formulas
  \item \textbf{Decomposes} revenue into volume + price; COGS into volume + rate
  \item \textbf{Generates} waterfall chart and writes CFO-ready memo in Word format
\end{enumerate}
\end{frame}

% ========================================
% SECTION 6: AUTOMATED REPORTING PIPELINE
% ========================================
\section{Automated Reporting Pipelines}

\begin{frame}{The Reporting Pipeline}
Most finance teams spend hours assembling weekly or monthly reports.  An AI-powered pipeline automates the entire chain.
\vspace{0.3cm}
\begin{center}
\begin{tikzpicture}[
  node distance=0.8cm,
  every node/.style={font=\small},
  sbox/.style={draw=titlegray, rounded corners, minimum width=2cm, minimum height=0.8cm, fill=excelinput, text=titlegray, font=\scriptsize\bfseries, align=center},
  sarrow/.style={-{Stealth[length=3mm]}, thick, draw=accentblue}
]
\node[sbox] (db) {Database\\Query};
\node[sbox, right=of db] (transform) {Transform\\\& Compute};
\node[sbox, right=of transform] (chart) {Charts \&\\Tables};
\node[sbox, right=of chart] (narr) {AI\\Narrative};
\node[sbox, right=of narr] (pptx) {PowerPoint\\Output};

\draw[sarrow] (db) -- (transform);
\draw[sarrow] (transform) -- (chart);
\draw[sarrow] (chart) -- (narr);
\draw[sarrow] (narr) -- (pptx);
\end{tikzpicture}
\end{center}
\vspace{0.5cm}
From raw data to polished deck --- \alert{no copy-paste, no formatting, no manual writing.}
\end{frame}

\begin{frame}{Building the Pipeline with Streamlit}
\textbf{Streamlit} turns a Python script into a web app in minutes.  Combined with an AI pipeline, it becomes a \alert{self-service reporting tool}.
\vspace{0.3cm}
\begin{columns}[T]
\begin{column}{0.5\textwidth}
\begin{shadedbox}[title=\textbf{User Interface}]
\begin{itemize}\small
  \item Dropdown: select time period, department, or metric
  \item Button: ``Generate Report''
  \item Download: auto-generated \texttt{.pptx} file
\end{itemize}
\end{shadedbox}
\end{column}
\begin{column}{0.5\textwidth}
\begin{shadedbox}[title=\textbf{Behind the Scenes}]
\begin{itemize}\small
  \item Query database for selected parameters
  \item Compute metrics and build charts
  \item Send chart data to LLM for narrative
  \item Assemble PowerPoint with \texttt{python-pptx}
\end{itemize}
\end{shadedbox}
\end{column}
\end{columns}
\vspace{0.3cm}
The user clicks one button and gets a polished deck.  \alert{No analyst needed.}
\end{frame}

\begin{frame}{Example: Monthly Ops Review Pipeline}
\begin{columns}[T]
\begin{column}{0.5\textwidth}
\textbf{What the Pipeline Produces}
\begin{enumerate}\small
  \item Title slide, KPI summary table, and trend charts
  \item Variance waterfall (budget vs.\ actual)
  \item AI-generated executive summary
\end{enumerate}
\end{column}
\begin{column}{0.5\textwidth}
\textbf{Time Comparison}
\begin{itemize}\small
  \item \textbf{Manual}: 6--8 hours per month (data pull, Excel, copy-paste into PowerPoint, write narrative)
  \item \textbf{Pipeline}: \alert{30 seconds} per run
  \item \textbf{ROI}: First month pays for the build time
\end{itemize}
\end{column}
\end{columns}
\end{frame}

\begin{frame}[fragile]{Automated PowerPoint with python-pptx}
\texttt{python-pptx} is a Python library for creating and editing PowerPoint files programmatically.
\vspace{0.3cm}
\begin{columns}[T]
\begin{column}{0.5\textwidth}
\begin{shadedbox}[title=\textbf{What It Can Do}]
\begin{itemize}\small
  \item Create slides from templates
  \item Insert charts, tables, and images
  \item Apply corporate formatting (fonts, colors, logos)
  \item Populate placeholders with live data
\end{itemize}
\end{shadedbox}
\end{column}
\begin{column}{0.5\textwidth}
\begin{shadedbox}[title=\textbf{The AI Advantage}]
\begin{itemize}\small
  \item AI writes the \texttt{python-pptx} code for you
  \item Tell Claude: ``Create a 5-slide deck from this data with a waterfall chart on slide 3''
  \item AI handles layout, formatting, and data binding
\end{itemize}
\end{shadedbox}
\end{column}
\end{columns}
\vspace{0.3cm}
\alert{You describe the deck; AI builds the automation.}  The pipeline runs unattended after that.
\end{frame}

% ========================================
% SECTION 7: PLUGINS AND SAASPOCALYPSE
% ========================================
\section{Plugins and the SaaSpocalypse}

\begin{frame}{What is a Plugin?}
A \textbf{plugin} is a packaged collection of skills, agents, hooks, and MCP servers that can be shared across teams.
\vspace{0.3cm}
\begin{columns}[T]
\begin{column}{0.5\textwidth}
\textbf{A Plugin Can Include}
\begin{itemize}\small
  \item Skills (\texttt{SKILL.md} files)
  \item Agents (\texttt{AGENT.md} files)
  \item MCP servers (external tools)
\end{itemize}
\end{column}
\begin{column}{0.5\textwidth}
\textbf{Why Use Plugins?}
\begin{itemize}\small
  \item Share skills with your team
  \item Namespace prevents conflicts
  \item Install from marketplaces
\end{itemize}
\end{column}
\end{columns}
\vspace{0.3cm}
\textbf{Practical workflow:} Start with a standalone skill for quick prototyping. When ready to share, wrap it in a plugin by adding a manifest file.
\end{frame}

\begin{frame}[shrink=10]{The ``SaaSpocalypse'': Market Reaction to Plugins}
On January 30, 2026, Anthropic released 11 open-source plugins for Claude.  Within days, \alert{\$285 billion} in market cap evaporated from software, legal tech, and financial services stocks.
\vspace{0.2cm}
\begin{columns}[T]
\begin{column}{0.5\textwidth}
\begin{shadedbox}[title=\textbf{Hardest-Hit Stocks}]
\scriptsize
\begin{tabular}{@{} l r @{}}
\toprule
\textbf{Company} & \textbf{Drop} \\
\midrule
LegalZoom & $-$20\% \\
Thomson Reuters & $-$16\% \\
RELX (LexisNexis) & $-$14\% \\
Intuit & $>-$10\% \\
Salesforce & $-$7\% \\
\bottomrule
\end{tabular}
\end{shadedbox}
\end{column}
\begin{column}{0.5\textwidth}
\begin{shadedbox}[title=\textbf{Why It Happened}]
\begin{itemize}\scriptsize
  \item Plugins showed AI could replicate \textit{core workflows} of specialized enterprise software
  \item \textbf{No proprietary code} --- just markdown instructions + database connectors
  \item Investors asked: ``Why pay \$50K/yr for software when a plugin does it?''
\end{itemize}
\end{shadedbox}
\end{column}
\end{columns}
\end{frame}

% ========================================
% SECTION 9: SUMMARY
% ========================================

\begin{frame}{Summary}
\begin{columns}[T]
\begin{column}{0.25\textwidth}
\begin{shadedbox}[title=\textbf{Dashboards}]
\begin{itemize}\scriptsize
  \item NL replaces fixed views
  \item AI writes queries
  \item Follow-ups are instant
\end{itemize}
\end{shadedbox}
\end{column}
\begin{column}{0.25\textwidth}
\begin{shadedbox}[title=\textbf{Skills}]
\begin{itemize}\scriptsize
  \item Encode expertise
  \item Consistent methodology
  \item Specialized agents
\end{itemize}
\end{shadedbox}
\end{column}
\begin{column}{0.25\textwidth}
\begin{shadedbox}[title=\textbf{Pipelines}]
\begin{itemize}\scriptsize
  \item DB $\rightarrow$ chart $\rightarrow$ narrative $\rightarrow$ deck
  \item Streamlit interface
  \item python-pptx output
\end{itemize}
\end{shadedbox}
\end{column}
\begin{column}{0.25\textwidth}
\begin{shadedbox}[title=\textbf{Plugins}]
\begin{itemize}\scriptsize
  \item Package skills
  \item \$285B impact
  \item Finance ground zero
\end{itemize}
\end{shadedbox}
\end{column}
\end{columns}
\vspace{0.5cm}
\begin{center}
\alert{The best dashboard is no dashboard --- it's a conversation with your data.}
\end{center}
\end{frame}

\end{document}
